\documentclass[answers]{exam}

\usepackage{../../../mypackages}
\usepackage{../../../macros}

% Pour le tableau
\usepackage{array}
\usepackage{longtable}
\usepackage{multirow}
\usepackage[table]{xcolor}

% -----------------------------
% Variables (valeurs par défaut — écrasées par gen_projet.py)
% -----------------------------
\newcommand{\varPrenom}{Sohantadeo}
\newcommand{\varNom}{Favereau}
\newcommand{\varClasse}{6ème}
\newcommand{\varAnticipationNote}{5}
\newcommand{\varAnticipationAppreciation}{Excellent sur l’anticipation. Tu commences le travail tôt et j’ai régulièrement des traces de ton avancée bien avant les deadlines. C’est exactement l’état d’esprit attendu dans ce projet : avancer progressivement plutôt que tout faire au dernier moment.}
\newcommand{\varInteractionNote}{4}
\newcommand{\varInteractionAppreciation}{Très bonnes interactions par email : tes messages sont clairs, avec du contexte et souvent des captures d’écran pour expliquer où tu bloques. Les questions deviennent de plus en plus précises. Tu pourrais en revanche venir un peu plus me voir pendant le cours pour échanger et voir ce que tu pourrais faire ensuite.}
\newcommand{\varAttitudeNote}{5}
\newcommand{\varAttitudeAppreciation}{Excellente attitude pendant les séances. Tu es volontaire, sérieux et prêt à aider tes camarades. On voit que tu t’impliques vraiment dans le projet. Essaie simplement de travailler parfois avec d’autres élèves pour mélanger un peu plus les groupes.}
\newcommand{\varFormulasNote}{3.5}
\newcommand{\varFormulasAppreciation}{Très bon travail dans l’ensemble et bonne compréhension des formules. Quelques erreurs d’étourderie subsistent : une addition incorrecte dans la conclusion de la partie objets et l’utilisation trop systématique de IFERROR. Attention aussi à ne pas écrire des nombres en dur dans les formules et à utiliser des fonctions comme AVERAGE.}
\newcommand{\varNoteFinale}{17.5}
\newcommand{\varAppreciationGenerale}{Très bon trimestre. Tu anticipes bien, les échanges par email sont efficaces et ton attitude en classe est très sérieuse. Il reste quelques petites erreurs techniques à corriger dans les formules, mais l’ensemble est solide. Continue comme ça.}

\title{Projets Scientifiques -- Compte rendu}
\author{Évaluation de \varPrenom{} \varNom{}}
\date{Mars 2026}

% Mise en forme des solutions en bleu
\SolutionEmphasis{\color{blue}}
\renewcommand{\solutiontitle}{\noindent}

\definecolor{lavande}{RGB}{180,190,255}
\definecolor{noteRed}{RGB}{220,50,50}
\definecolor{notePaleRed}{RGB}{255,180,180}
\definecolor{noteYellow}{RGB}{255,230,100}
\definecolor{notePaleGreen}{RGB}{180,230,180}
\definecolor{noteDarkGreen}{RGB}{40,140,40}

% Colonne p{} alignée à gauche et qui wrappe
\newcolumntype{L}[1]{>{\raggedright\arraybackslash}p{#1}}

% -----------------------------
% Commande de coloration conditionnelle des notes
% -----------------------------
\newcommand{\noteCell}[1]{%
  \edef\argval{#1}%
  \def\defaultval{--}%
  \ifx\argval\defaultval
    #1%
  \else
    \pgfmathtruncatemacro{\noteLevel}{ceil(\argval)}%
    \ifcase\noteLevel
      % 0
      \cellcolor{noteRed}\textcolor{white}{\textbf{#1}}%
    \or
      % 0.5–1
      \cellcolor{noteRed}\textcolor{white}{\textbf{#1}}%
    \or
      % 1.5–2
      \cellcolor{notePaleRed}\textbf{#1}%
    \or
      % 2.5–3
      \cellcolor{noteYellow}\textbf{#1}%
    \or
      % 3.5–4
      \cellcolor{notePaleGreen}\textbf{#1}%
    \or
      % 4.5–5
      \cellcolor{noteDarkGreen}\textcolor{white}{\textbf{#1}}%
    \fi
  \fi
}

\begin{document}

\textbf{Collège Lycée Suger}
\hfill
\textbf{Projets Scientifiques} \\

\textbf{Année 2025-2026}
\hfill
\textbf{Classe de \varClasse{}} \par

{\let\newpage\relax\maketitle}

\begin{tcolorbox}[
  colframe=red!90!black,
  fonttitle=\bfseries,
  boxsep=4pt,
  sharp corners,
]
\begin{center}
\textbf{\textcolor{red}{Ce document doit être imprimé et collé dans le cahier de projets scientifiques.}}
\end{center}
\end{tcolorbox}

\subsection*{Note}

\renewcommand{\arraystretch}{1.7}

\noindent
\begin{longtable}{|L{0.14\textwidth}|L{0.57\textwidth}|>{\centering\arraybackslash}p{0.08\textwidth}|>{\centering\arraybackslash}p{0.1\textwidth}|}
\hline
\rowcolor{lavande}
\textbf{Catégorie} & \textbf{Appréciation} & \textbf{Note} & \textbf{Barème} \\[2pt]
\hline

Anticipation
  & \varAnticipationAppreciation{}
  & \noteCell{\varAnticipationNote}
  & 5 \\
\hline

Interaction avec l'enseignant
  & \varInteractionAppreciation{}
  & \noteCell{\varInteractionNote}
  & 5 \\
\hline

Attitude et effort
  & \varAttitudeAppreciation{}
  & \noteCell{\varAttitudeNote}
  & 5 \\
\hline

Formules et sources
  & \varFormulasAppreciation{}
  & \noteCell{\varFormulasNote}
  & 5 \\
\hline

% ---------------- Total ----------------
\rowcolor{green!35}
\textbf{Total} & & \textbf{\varNoteFinale{}} & \textbf{20} \\

\hline

\end{longtable}

\subsection*{Appréciation générale}

\begin{tcolorbox}[
  colback=blue!5!white,
  colframe=blue!70!black,
  title=Commentaire,
  fonttitle=\bfseries,
  boxsep=4pt,
  arc=4pt,
]
\varAppreciationGenerale{}
\end{tcolorbox}


\subsection*{Rappel -- Barème de notation}

\renewcommand{\arraystretch}{1.7}

\noindent
\begin{longtable}{|L{0.14\textwidth}|L{0.73\textwidth}|>{\centering\arraybackslash}p{0.04\textwidth}|}
\hline
\rowcolor{lavande}
\textbf{Catégorie} & \textbf{Description} & \textbf{/5} \\[2pt]
\hline

Anticipation
  & L'élève a-t-il travaillé en avance et planifié son travail ? On attend une progression régulière visible dans l'historique Google Sheets, pas un travail fait au dernier moment. Un travail clairement anticipé avec planification visible donne 5/5.
  & 5 \\
\hline

Interaction avec l'enseignant
  & Qualité des échanges par email et en classe. On attend des messages clairs, contextualisés, avec des captures d'écran si nécessaire et une explication de ce qui a été essayé. Les retours de l'enseignant doivent être pris en compte.
  & 5 \\
\hline

Attitude et effort
  & Implication générale pendant le projet : autonomie, persévérance, aide apportée aux camarades, capacité à solliciter de l'aide au bon moment. Un élève qui demande de l'aide intelligemment n'est pas pénalisé.
  & 5 \\
\hline

Formules et sources
  & Maîtrise des formules Google Sheets (VLOOKUP, SOMME, multiplication, AVERAGE\ldots), compréhension de la logique derrière les calculs et qualité des sources utilisées. Éviter les valeurs en dur et les IFERROR non justifiés.
  & 5 \\
\hline

% ---------------- Total ----------------
\rowcolor{green!35}
\textbf{Total} & \textbf{Note maximale possible} & \textbf{20} \\
\hline

\end{longtable}


\end{document}
